% !TeX root = ../main.tex
% Add the above to each chapter to make compiling the PDF easier in some editors.
% Read the CSV files

\chapter{Introduction}\label{chapter:introduction}

Just-in-time (JIT) compilation is a widely used technique for improving performance for applications, where the combined speed of compilation and execution is critical. Such
applications include dynamic runtimes such as the Java Virtual Machine~\cite{jvm} as well as
the acceleration of database queries~\cite{sql-jit}. 

Since compile time directly contributes to end-to-end latency, fast machine code generation is especially important for JIT compilation. Recently, TPDE~\cite{tpde} has shown that the compilation speed of the LLVM~\cite{LLVM} unoptimized back-end can be improved substantially, while keeping a similar runtime performance. 

Since JIT systems inherently involve a trade-off between compilation and runtime, many
systems employ multi-stage approaches, with the first stage serving as the baseline and subsequent stages being progressively optimized~\cite{tpde}.
However, when using LLVM, the difference in compilation time between TPDE and LLVM's optimized backend (\texttt{-O1}) is two orders of magnitude, leaving a significant gap for programs that could benefit from better code generation, but for which the substantial increase in compilation time cannot be justified.

To address this gap, we propose a novel register allocation approach built on top of TPDE that enables improved code generation while maintaining relatively fast compilation times. Our approach consists of two analysis passes: a liveness analysis, which additionally computes next-use information, and a spill analysis pass that determines which values to spill and optimizes spill locations to reduce spill overhead in loops. During code generation, we introduce a repair mechanism for handling instruction constraints, enabling more efficient code generation for constrained operations such as integer division on \texttt{x86-64}.

On SPECint 2017, we achieve a 24\% runtime improvement over the LLVM-O0 backend,
while the compilation time for optimized LLVM-IR remains three times faster. However,
this also results in a 5x increase in compilation time compared to TPDE. Nevertheless, our work achieves a
new Pareto optimum for the compilation of optimized IR, as it is significantly faster than LLVM-O1 in terms of compilation time,
while simultaneously offering a runtime improvement over both
TPDE and LLVM-O0.
\newpage
The key contributions of our work include:

\begin{itemize}
    \item A lifetime analysis that generates next-use information in the same pass
    \item A spill analysis that identifies values to spill without relying on knowledge of the instruction semantics, or requiring modification to the IR.
    \item A simplified repairing mechanism to handle instruction constraints during code generation, without rolling back code generation.
\end{itemize}

The remainder of this work is structured as follows: In \autoref{ch:background} we provide background on register allocation and TPDE, as well as the constraints arising from the framework’s design. In \autoref{chapter:approach} we present our approach. In \autoref{chapter:related_work} we discuss related work in register allocation, and how it relates to our approach. Next we evaluate our approach compared to TPDE, LLVM-O0 and LLVM-O1 in \autoref{chapter:evaluation}. Finally, we summarize our results and discuss future work in \autoref{chapter:conclusion}.